% spec body template — the slug-echo body filename (<slug>.tex, per #295)
% lives at <thread>.{N}/<slug>.tex. This is a normative specification: every
% claim is true of the implementation (code_ref) or explicitly marked
% target-state. A spec's body is LaTeX (SKILL.md §Output format) — multi-file
% friendly via \input{sections/*.tex}. Delete these comments before authoring.
% Replace bracketed placeholders.
%
% ADOPTION NOTE: if you are adopting an existing multi-file LaTeX spec, you do
% NOT start from this template — place the existing tree per SKILL.md
% §"Adopting an existing spec". This skeleton is only for the deferred
% draft-from-scratch / stub-scaffold case (spec-draft step 4 "Scaffold").

\documentclass[11pt]{article}

\usepackage[margin=1in]{geometry}
\usepackage{graphicx}
\usepackage{amsmath,amssymb}
\usepackage{booktabs}
\usepackage{hyperref}

\title{[System] Specification}
\author{[Maintainer / team]}
\date{\today}

\begin{document}
\maketitle
\tableofcontents

\section{Scope and conformance}
% What this spec normatively governs. State the RFC 2119 keyword discipline
% (MUST / SHALL / SHOULD / MAY) once here and use it consistently. Name the
% implementation this spec describes (the code_ref companion) so a reader
% knows the source of truth the normative claims are checked against.

\section{Definitions and notation}
% Define EVERY term used in a normative predicate below. An undefined term
% used in a MUST/SHALL is the "Undefined normative term" critical flag.

\section{[Component / message / structure]}
% The normative content. For each constant, struct/field layout, formula, and
% validity predicate: state it unambiguously and falsifiably (dim 3), and keep
% it consistent with every other section that states the same quantity (dim 2).
% Each claim either matches the implementation (code_ref) or is explicitly
% marked target-state with the gap tracked.

\subsection{Constants}
% A grep-able constants table makes internal-consistency and (Phase 3)
% deterministic constant-extraction checks cheap. State each constant once,
% authoritatively; reference it elsewhere rather than restating a literal.
%
% Annotate the AUTHORITATIVE statement of each normative constant with an
% "% anvil-const:" marker so spec-review's deterministic gate can catch the
% same named constant stated with a different value in another section (the
% block-time-floor 3s-vs-5s / ring-size-byte-count drift the gate exists to
% prevent). The marker works on its own line OR as a trailing table-row
% comment (everything before the "%" is ignored):
%
%   % anvil-const: name=block_time_floor value=3 unit=s
%
% Rules (see SKILL.md "Constant-consistency markers"): name= and value= are
% required; unit= is optional; comparison is exact after minimal normalization
% (no numeric tolerance); same name + different unit is a separate lower-
% severity finding; a \newcommand duplicate/conflicting redefinition is caught
% too. Prose-level (unmarked) drift is NOT caught by the gate — it is left to
% review judgment and spec-audit.
\begin{center}
\begin{tabular}{@{}lll@{}}
\toprule
Name & Value & Notes \\
\midrule
[CONST\_A] & [value] & [what it governs] \\ % anvil-const: name=[CONST_A] value=[value] unit=[unit]
\bottomrule
\end{tabular}
\end{center}

\section{Validity predicates}
% The formal block/message/state-validity predicates. Every symbol used here
% must be defined in §Definitions. A predicate must be satisfiable (dim 5).

\section{Implementation status}
\label{sec:impl-status}
% THE IMPLEMENTATION-STATUS REGISTER (SKILL.md §Implementation-status register).
% Live vs. target per component. Where the implementation (code_ref) and this
% spec deliberately diverge, record the gap HERE rather than silently rewriting
% either side. This table is OPERATOR/DRAFTER-authored input, NOT
% auditor-generated: you declare which components have an intentional live/target
% gap; spec-audit checks that a target-state claim in the body has a matching row
% (a code-vs-spec divergence with NO row is the "unregistered gap" critical flag,
% Disposition intentional-gap). Only components with a REAL live/target
% divergence need a row — a true-now claim needs none. Every target-state row
% names the divergence explicitly (concrete Live + concrete Target, no vague
% "TBD"); the Tracking column links the issue/ADR that owns the gap's resolution.
% Delete the example row and add your own; delete the whole section only if the
% spec claims no live/target gaps at all.
\begin{center}
\begin{tabular}{@{}p{2.2cm}p{2.8cm}p{2.8cm}ll@{}}
\toprule
Component & Live behavior & Target (this spec) & Status & Tracking \\
\midrule
[Output signatures] & [Ed25519 (classical)] & [ML-DSA-65 on all outputs] & target-state & [issue link / unscheduled] \\
\bottomrule
\end{tabular}
\end{center}

\section{Revision history}
% Versioning discipline (dim 7): a reader must know which version they hold and
% what changed between versions.
\begin{center}
\begin{tabular}{@{}lll@{}}
\toprule
Version & Date & Summary of normative changes \\
\midrule
1 & \today & Initial draft \\
\bottomrule
\end{tabular}
\end{center}

\end{document}
